\documentclass[acmsmall,screen,review,anonymous]{acmart}

\setcopyright{none}
\acmConference[POPL '27]{ACM SIGPLAN Symposium on Principles of Programming Languages}{January 2027}{}
\acmYear{2027}
\acmDOI{}
\acmISBN{}

\usepackage{amsmath}
\usepackage{amssymb}
\usepackage{mathtools}
\usepackage{booktabs}
\usepackage{enumitem}
\usepackage{xspace}

\newcommand{\system}{CiPR-FTQC\xspace}
\newcommand{\lang}{FTQP\xspace}
\newcommand{\judgment}{\mathcal{K};\Gamma;\Delta \vdash t : \Gamma';\Delta' \;!\; E}
\newcommand{\code}[1]{\ensuremath{\mathsf{#1}}}
\newcommand{\gate}[1]{\ensuremath{\mathsf{#1}}}
\newcommand{\res}[1]{\ensuremath{\mathsf{#1}}}
\newcommand{\mode}[1]{\ensuremath{\mathsf{#1}}}
\newcommand{\cert}[1]{\ensuremath{\mathsf{#1}}}
\newcommand{\Cap}{\ensuremath{\mathsf{Cap}}}
\newcommand{\Switch}{\ensuremath{\mathsf{Switch}}}
\newcommand{\Factory}{\ensuremath{\mathsf{Factory}}}
\newcommand{\Backend}{\ensuremath{\mathsf{Backend}}}
\newcommand{\Checked}{\cert{Checked}}
\newcommand{\Certified}{\cert{Certified}}
\newcommand{\Assumed}{\cert{Assumed}}
\newcommand{\State}{\mode{State}}
\newcommand{\Observable}{\mode{Observable}}

\title[Code-Indexed Resource Calculus for FTQC]{A Code-Indexed Probabilistic Resource Calculus for Fault-Tolerant Quantum Programming}

\author{Anonymous Authors}
\affiliation{\institution{Anonymous Institution}\country{}}

\begin{document}

\begin{abstract}
Fault-tolerant quantum programs are written as logical operations, but they run through a heterogeneous collection of code-specific gates, code switches, magic-state factories, decoders, and layout constraints.
This mismatch makes a compiler decision hard to audit: a direct gate may be applied outside the current code's capability set, a probabilistic resource state may be reused as if it were classical data, or an imported physical theorem may be silently treated as a local proof.
This paper presents \system, a code-indexed, probabilistic, resource-aware calculus for elaborating logical quantum programs into certified fault-tolerant acquisition plans.
The calculus follows a simple running example: a \code{SurfaceD5} program that asks for $\gate{T}$, $\gate{CNOT}$, and $\gate{CCZ}$ operations even though its initial code exposes only a small direct gate set.
\system assigns encoded data types $Q[C,d,\mu]$, tracks magic and Bell resources linearly, and composes conservative effects for logical error, failure, acceptance, space, and time.
Its central judgment, $\judgment$, checks that each target acquisition step is justified by a capability rule, consumes resources exactly once, respects a fixed backend context, and records whether the rule is locally checked, certified by an imported theorem, or still assumed.
We implement the calculus in a prototype compiler/checker that plans direct gates, switch-via-code paths, resource-state injection, and hybrid regions over a small rule library drawn from recent fault-tolerant quantum protocols.
The implementation emits auditable sidecars, SMT obligations, protocol certificates, decoder contracts, geometry checks, and a Lean proof kernel for core resource and certificate properties.
On the current case studies, the checker accepts multiple exploratory acquisition strategies, rejects unsupported direct gates, duplicate resource consumption, and backend-incompatible qLDPC-style switches, and produces a full-stack report with 53 proof obligations and explicit imported-theorem boundaries.
\end{abstract}

\maketitle

\section{Introduction}

Fault-tolerant quantum computation forces a compiler to reason about two languages at once.
At the algorithmic level, a programmer writes logical operations such as $\gate{H}$, $\gate{CNOT}$, $\gate{T}$, measurements, error-correction rounds, and classical control.
At the implementation level, those operations are available only through encoded data and code-specific mechanisms: a surface-code patch may expose a convenient Clifford subset, a Steane-like code may expose transversal Clifford gates, a tetrahedral color-code instance may expose a transversal non-Clifford gate, and a magic-state factory may implement a missing operation by preparing and consuming a noisy ancillary resource.
The compiler must therefore decide not merely where to place gates, but how to acquire the fault-tolerant capability that makes each logical gate legal.

The difficulty is easiest to see through a small program that starts all live data in a \code{SurfaceD5} profile.
The program prepares four logical qubits, performs a surface-code error-correction round, applies a direct $\gate{H}$, then enters a dense region containing $\gate{CNOT}$ and $\gate{T}$ gates before a later $\gate{CCZ}$ and a measurement-controlled correction.
In the rule library used by our prototype, \code{SurfaceD5} supports direct $\gate{H}$ and $\gate{S}$, but not direct $\gate{CNOT}$, $\gate{T}$, or $\gate{CCZ}$.
Thus the source program is not wrong, but it is incomplete as an executable fault-tolerant program: the compiler must elaborate each missing operation into an acquisition plan such as resource-state injection, code switching, or a hybrid region that switches once and executes several gates before switching back.

This acquisition step is a programming-language problem because its correctness depends on information that is usually hidden in compiler engineering artifacts.
If a compiler emits a single opaque ``apply this gadget'' step, a checker cannot tell whether the current code actually supports the requested gate, whether a magic state was consumed exactly once, whether a code switch fits the target backend, or whether a proof obligation was discharged locally or imported from a protocol paper.
The issue is not that fault-tolerant protocols lack proofs; recent protocol papers provide many useful building blocks.
The issue is that a compiler needs a small common interface for composing those building blocks without erasing their capability requirements, resource behavior, effect bounds, and proof status.

This paper presents \system, a code-indexed probabilistic resource calculus for fault-tolerant quantum programming.
The source language expresses logical intent, while the target language makes acquisition explicit through atoms for native or transversal gates, code switching, resource production and consumption, logical measurement, decoding, frame update, postselection, and reset.
The central type of an encoded value is $Q[C,d,\mu]$: the qubit is currently protected by code $C$, at distance parameter $d$, with correctness mode $\mu$ distinguishing state-level from observable-level reasoning.
The central judgment, $\judgment$, says that a target program transforms quantum context $\Gamma$ and linear resource context $\Delta$ into new contexts while producing a conservative effect summary $E$.
In the running example, this judgment is what prevents a direct $\gate{CNOT}$ on \code{SurfaceD5}, while allowing a checked plan that first acquires a Bell resource or switches into a code with the right transversal capability.

The key design choice in \system is that proof status is part of the rule interface.
A rule may be \Checked, when the repository discharges the relevant algebraic or software-level obligation locally; \Certified, when local checks are combined with an explicit imported theorem from a cited protocol; or \Assumed, when the rule is a prototype bridge that should not be counted as a paper-level physics claim.
The strict policy rejects \Assumed{} rules, while exploratory compilation may retain them with warnings.
This policy is deliberately prosaic but important: it turns the phrase ``assuming this protocol is valid'' into a machine-visible certificate boundary that follows the compiled plan.

We have implemented \system as \lang, a small compiler/checker and formal toolchain.
The planner elaborates logical programs into direct, resource, switch, hybrid, and randomized acquisition strategies.
The checker validates capability witnesses, certificate policies, resource linearity, fixed-backend topology constraints, layout-event accounting, and effect bounds.
The toolchain emits JSON sidecars, SMT-LIB obligations, GF(2) protocol certificates, decoder reports, rectangular patch-geometry checks, artifact hashes, and a Lean 4 proof kernel for core calculus properties.
The current physical boundary is explicit: algebraic code-switching and resource-linearity properties are checked locally, while complete flag-circuit fault enumeration, surface-code decoder internals, and architecture-level factory profiles remain imported theorem dependencies.

The case studies are intended to support these language and checker claims, not to announce a new fault-tolerant physical protocol.
On the small running workload, \system validates magic-only, switch-only, hybrid, best, and randomized acquisition strategies on a fixed $64\times64$ grid backend.
The hybrid strategy uses two switches to reduce the cycle count of the dense region, while the magic-only strategy avoids switches at the cost of more factories.
Negative tests reject exactly the mistakes the calculus is designed to expose: a direct $\gate{CNOT}$ on \code{SurfaceD5}, duplicate consumption of a shared $\res{TState}$, and a qLDPC-style switch on a grid backend.

This paper makes the following contributions.
\begin{itemize}[leftmargin=*]
  \item We formulate fault-tolerant acquisition as a code-indexed target calculus whose primitive atoms expose the evidence needed to audit code capabilities, code switching, resource-state injection, measurement, decoder, frame, postselection, and reset operations.
  \item We give a type-and-effect discipline for acquisition plans, using code-indexed quantum types, linear probabilistic resources, fixed backend capabilities, certificate policies, and conservative effect composition.
  \item We mechanize a small proof kernel in Lean for capability rejection, strict-certificate rejection, resource linearity, and effect-bound composition, while marking physical protocol claims that remain imported theorems.
  \item We implement a prototype compiler/checker and formal toolchain that connects rule-library metadata, protocol certificates, decoder contracts, geometry checks, SMT obligations, and reproducible artifact hashes.
  \item We report case studies and negative tests showing how the checker accepts auditable acquisition strategies and rejects illegal direct gates, duplicated resources, and backend-incompatible switches.
\end{itemize}

\section{Overview}

This section follows one program through the \system pipeline.
The program is small enough to inspect, but it contains the three phenomena that make fault-tolerant acquisition interesting: a direct gate, a missing non-Clifford gate that can be supplied by a resource, and a dense region where switching once may be better than acquiring each gate independently.
In pseudocode, the source program prepares four \code{SurfaceD5} logical qubits, runs an error-correction round, applies a direct $\gate{H}$, enters a region with two $\gate{CNOT}$ gates and two $\gate{T}$ gates, later applies a $\gate{CCZ}$, and finishes with measurement-controlled $\gate{S}$ correction.
The source says what logical computation is intended; it does not say how to realize the missing operations.

Figure~\ref{fig:pipeline} shows the compilation path.
The planner takes the logical program and a backend context, consults a rule library, and emits an acquisition plan with a sidecar.
The sidecar is not an afterthought: it is the evidence object that tells the checker which rule justified each step, which resource was produced or consumed, how effects were composed, which layout event was reported, and which certificate level was used.
The formal toolchain then checks the sidecar through independent components: Python validates step-level legality, Z3 checks arithmetic obligations, GF(2) checkers discharge algebraic protocol facts, decoder checkers validate local decoder contracts where possible, geometry checkers verify conservative rectangular embeddings, and Lean proves small core properties about the calculus.

\begin{figure}[t]
\centering
\fbox{\begin{minipage}{0.93\linewidth}
\centering
\[
\begin{array}{c}
\text{logical program over \code{SurfaceD5} qubits} \\
\Downarrow \\
\text{planner chooses direct, resource, switch, or hybrid acquisition} \\
\Downarrow \\
\text{typed acquisition plan + effect/layout/certificate sidecar} \\
\Downarrow \\
\text{checker + SMT + GF(2) certificates + decoder + geometry + Lean}
\end{array}
\]
\end{minipage}}
\caption{The \system pipeline. The running example starts as a logical surface-code program; the planner makes missing capabilities explicit; the checker audits the evidence that justifies each acquisition step.}
\label{fig:pipeline}
\end{figure}

The direct gate illustrates what the checker should allow without ceremony.
When the source applies $\gate{H}$ to a \code{SurfaceD5} qubit, the rule library contains a direct surface-code capability for $\gate{H}$.
The target plan therefore contains a single direct step, and the type of the qubit remains $Q[\code{SurfaceD5},5,\State]$.
No resource is produced, no switch is required, and the effect contribution is just the rule's local cost.
This mundane case matters because it fixes the baseline: a direct step is legal only when the current code index supports it.

The $\gate{T}$ gate illustrates why resources must be linear.
Because \code{SurfaceD5} does not expose direct $\gate{T}$, one acquisition candidate is:
\[
  \mathsf{produce}\ r:\res{TState} ;\quad
  \mathsf{consume}\ r\ \mathsf{for}\ \gate{T}\ q.
\]
The notation is intentionally close to the implementation sidecar.
The first step creates a fresh physical artifact; the second step spends that artifact to implement the logical gate.
If a later target step tries to consume the same identifier again, the resource context no longer contains it, and the checker rejects the plan.
This is the simplest example where a formula in the paper corresponds directly to a bug that the prototype catches.

The dense region illustrates why the compiler should not always acquire gates one by one.
For a block such as
\[
  \gate{CNOT}(q_0,q_1);\ \gate{T}(q_1);\ \gate{CNOT}(q_1,q_2);\ \gate{T}(q_2),
\]
the hybrid strategy may switch the involved qubits into \code{Tetra15}, execute the region using transversal capabilities, and switch them back:
\[
  \mathsf{switch}\ \vec q:\code{SurfaceD5}\to\code{Tetra15};\quad
  \mathsf{transv}\ \text{region};\quad
  \mathsf{switch}\ \vec q:\code{Tetra15}\to\code{SurfaceD5}.
\]
The checker sees this not as a magical optimization, but as a sequence of ordinary typed steps.
During the middle of the plan the qubits have code index \code{Tetra15}; after the final switch their code index returns to \code{SurfaceD5}.
The case study in Section~\ref{sec:case-study} shows how this choice changes cycle counts, factory counts, switch counts, and assumption warnings.

The same example also explains the role of certificate levels.
The locally checked part of a switch certificate may verify an algebraic gauge-fixing core over GF(2), while the complete flag-circuit single-fault-tolerance theorem remains imported from the cited protocol paper.
Other rules in the current prototype are weaker: they are \Assumed{} bridge rules used to connect otherwise paper-backed components in a compact compiler experiment.
An exploratory run may keep these rules so that the planner can compare acquisition strategies, but a strict run rejects them.
Thus the overview example already contains the paper's main claim: \system does not make every physical protocol true; it makes the proof boundary visible when protocols are composed.

\section{Core Language}

\subsection{Codes, Capabilities, and Certificates}

\system separates the code a value currently inhabits from the rules that can act on that code.
In the running example, the four prepared data qubits begin with type $Q[\code{SurfaceD5},5,\State]$.
This type is a compact way to say three things the checker needs later: the qubit is encoded in the surface profile, the rule library treats it as distance five, and later operations may rely on state-level correctness rather than only final observable correctness.
The type is not a physical simulator state; it is the static handle that lets the compiler ask whether a requested operation is a direct capability, a switch target, or a resource-acquisition goal.

A code specification records the information needed to answer those questions.
The current prototype contains \code{SurfaceD5}, \code{Steane3}, \code{Tetra15}, and a negative-control \code{QLDPC12} profile.
\code{SurfaceD5} supports direct $\gate{H}$ and $\gate{S}$ in the small rule library; \code{Steane3} supports transversal Clifford gates; \code{Tetra15} supports transversal $\gate{T}$ and $\gate{CNOT}$; and \code{QLDPC12} requires long-range topology.
The last profile is deliberately included so that the checker has a concrete backend-incompatibility case to reject.

The capability context $\mathcal{K}$ collects the rule facts available to a compilation run:
\[
\begin{array}{rcl}
\mathcal{K} &::=& \Cap(C,g) \mid \Switch(C_1,C_2) \mid \Factory(R) \mid \Backend(B) \mid \cdots
\end{array}
\]
The interesting point is that these facts are not bare booleans.
Each fact carries a certificate level $\ell\in\{\Checked,\Certified,\Assumed\}$.
\Checked{} rules are discharged by local checkers or simple implementation invariants.
\Certified{} rules combine local checks with an explicit imported theorem from a cited protocol.
\Assumed{} rules are prototype glue and are rejected by the strict theorem policy.
This is the mechanism that prevents the paper from confusing a useful compiler experiment with a new physical proof.

\subsection{Source and Target Syntax}

The source language is intentionally ordinary because users should describe logical intent rather than acquisition mechanics.
It includes preparation, logical gates, error-correction rounds, logical measurement, sequencing, and classical branching:
\[
\begin{array}{rcl}
e &::=& \mathsf{prepare}\ q\ s
  \mid \mathsf{gate}\ g\ \vec q
  \mid \mathsf{ec}\ \vec q\ d
  \mid \mathsf{measure}\ b\ O\ \vec q \\
  && \mid\ \mathsf{if}\ b\ \mathsf{then}\ e_1\ \mathsf{else}\ e_2
  \mid e_1;e_2.
\end{array}
\]
In the case study, this language is represented by \texttt{LogicalOp.prepare}, \texttt{LogicalOp.apply}, \texttt{LogicalOp.ec}, \texttt{LogicalOp.measure}, and \texttt{LogicalOp.branch}.
The source contains $\gate{T}$ and $\gate{CNOT}$ as logical gates even though the initial code does not directly support them.
That is the point of elaboration: unsupported source gates become acquisition obligations rather than immediate type errors.

The target language is smaller and more explicit.
It exposes the atoms that a checker can audit:
\[
\begin{array}{rcl}
a &::=& \mathsf{alloc}\ q:C
  \mid \mathsf{native}\ g\ \vec q
  \mid \mathsf{transv}\ g\ \vec q
  \mid \mathsf{switch}\ q:C_1\to C_2 \\
  && \mid\ \mathsf{produce}\ r:R
  \mid \mathsf{consume}\ r\ \mathsf{for}\ g\ \vec q
  \mid \mathsf{measureL}\ b\ O\ \vec q \\
  && \mid\ \mathsf{decode}\ s
  \mid \mathsf{frame}\ f
  \mid \mathsf{postselect}\ p
  \mid \mathsf{reset\_clean}\ q, \\
t &::=& a \mid t_1;t_2 \mid \mathsf{if}\ b\ \mathsf{then}\ t_1\ \mathsf{else}\ t_2.
\end{array}
\]
The target has no unrestricted $\mathsf{apply}\ U$ primitive.
That omission is a design constraint, not an aesthetic preference.
If arbitrary unitaries could appear as primitive atoms, a compiler could hide an entire code-switching or factory protocol behind one step, and the checker would lose the ability to inspect capability, resource, effect, and certificate evidence.

\subsection{Types, Resources, and Effects}

The encoded quantum type $Q[C,d,\mu]$ is the thread that connects the paper calculus to the implementation.
The code index $C$ changes when the planner inserts a switch, stays fixed for direct and resource-injection gates, and must agree across branches before the program continues.
The mode $\mu ::= \State \mid \Observable$ distinguishes state-level values from values whose guarantee is only about final measurement distributions.
The implementation already carries this field in \texttt{QType}; the calculus uses it to state the intended discipline that observable-level outputs should not be fed to rules requiring state-level inputs.

Probabilistic resources live in a separate context because they are physical artifacts, not reusable descriptions:
\[
R ::= \res{TState}[\epsilon] \mid \res{CCZState}[\epsilon] \mid \res{BellPair}[\epsilon] \mid \res{AuxState}[C,\epsilon].
\]
In the implementation, a resource acquisition step places a fresh identifier such as \texttt{res\_t\_0} in the sidecar, and the corresponding consume step names the same identifier.
The checker maintains produced and consumed sets while flattening the plan.
This small piece of bookkeeping is the operational version of the linear resource context $\Delta$.

Effects summarize what a plan promises conservatively:
\[
E = \langle \epsilon, f, a, s, \tau, \rho, n_{\mathsf{sw}}, n_{\mathsf{fac}}, \ldots\rangle.
\]
Here $\epsilon$ is a logical-error upper bound, $f$ is a failure upper bound, $a$ is an acceptance lower bound, $s$ is peak space, $\tau$ is cycle count, $\rho$ is qubit-round count, and the remaining fields count switches, factories, measurements, resets, multi-qubit gates, decoder latency, certificates, assumptions, and rule identifiers.
Sequential composition adds error, cycles, qubit-rounds, and counts, multiplies acceptance lower bounds, and takes a maximum for peak space.
Branch composition takes conservative branchwise bounds.
The effect algebra is deliberately simple because it must be checked on every generated sidecar; its value is that the numbers in Section~\ref{sec:case-study} are produced by the same composition discipline used by the checker.

\section{Typing and Metatheory}

\subsection{The Audit Judgment}

The core target judgment is:
\[
  \judgment.
\]
This judgment is best read as an audit contract for the generated sidecar.
Under capability context $\mathcal{K}$, quantum context $\Gamma$, and linear resource context $\Delta$, target program $t$ transforms the contexts to $\Gamma'$ and $\Delta'$ and reports effect $E$.
For the running example, the judgment says that the checker can follow each qubit as it leaves \code{SurfaceD5}, enters \code{Tetra15} for a hybrid region, returns to \code{SurfaceD5}, and consumes resource identifiers exactly once along the way.
The judgment is parameterized by a certificate policy $\pi$, so the same syntactic plan may be accepted in exploratory mode and rejected in strict mode if it uses \Assumed{} rules.

Figure~\ref{fig:typing} shows the rules that matter most for the story.
The direct rule captures why $\gate{H}$ on \code{SurfaceD5} is routine but direct $\gate{CNOT}$ is rejected: the current code index must provide a capability witness.
The switch rule captures why a dense region can temporarily use another code: the type of the qubit changes from $Q[C_1,d,\State]$ to $Q[C_2,d',\State]$.
The produce and consume rules capture why a magic state behaves linearly: once the identifier leaves $\Delta$, a second consume step cannot typecheck.

\begin{figure}[t]
\[
\frac{
  q:Q[C,d,\mu]\in\Gamma \quad \Cap(C,g)\in\mathcal{K} \quad \pi(\ell)=\mathsf{true}
}{
  \mathcal{K};\Gamma;\Delta \vdash \mathsf{direct}\ g\ q : \Gamma;\Delta \;!\; E_g
}
\]
\[
\frac{
  q:Q[C_1,d,\State]\in\Gamma \quad \Switch(C_1,C_2)\in\mathcal{K} \quad \pi(\ell)=\mathsf{true}
}{
  \mathcal{K};\Gamma;\Delta \vdash \mathsf{switch}\ q:C_1\to C_2 :
  \Gamma[q\mapsto Q[C_2,d',\State]];\Delta \;!\; E_{\mathsf{sw}}
}
\]
\[
\frac{
  r\notin\mathrm{dom}(\Delta) \quad \Factory(R)\in\mathcal{K}
}{
  \mathcal{K};\Gamma;\Delta \vdash \mathsf{produce}\ r:R :
  \Gamma;\Delta,r:R \;!\; E_{\mathsf{prod}}
}
\qquad
\frac{
  r:R\in\Delta
}{
  \mathcal{K};\Gamma;\Delta \vdash \mathsf{consume}\ r\ \mathsf{for}\ g\ q :
  \Gamma;\Delta-r \;!\; E_{\mathsf{use}}
}
\]
\caption{Representative typing rules for capability use, code switching, and linear probabilistic resources. The rules are chosen because each corresponds to a concrete checker behavior in the case study.}
\label{fig:typing}
\end{figure}

\subsection{What Is Mechanized Today}

The Lean development mechanizes the small core needed to make these checker invariants precise.
It defines codes, gates, modes, resource kinds, certificate levels, policies, primitive atoms, resource contexts, effects, and typed programs.
The most important theorem for the running example is the unsupported-direct-gate rejection theorem: if the capability predicate for $(C,g)$ is false, there is no well-typed direct step for $g$ on code $C$.
This theorem is the formal counterpart of the negative test that rejects direct $\gate{CNOT}$ on \code{SurfaceD5}.

The second important theorem family concerns resource linearity.
Lean proves that after a resource has been consumed, consuming the same identifier again is impossible in the resulting resource context.
The implementation mirrors this theorem with produced and consumed sets in the checker.
The point is not that the proof is mathematically deep; the point is that a common compiler error has a direct expression in the calculus, the checker, and the test suite.

The third theorem family concerns certificates and effects.
Lean proves that typed programs contain only certificate levels allowed by the selected policy, and that sequential and branch effect composition conservatively bounds component effects.
For the case study, this means a strict run rejects plans that depend on prototype bridge rules, and an exploratory run can still report the same plan with visible warnings.
It also means the total effect summary is not a detached table entry; it is the result of applying the same composition operators used in the formal core.

\subsection{Conditional Soundness Boundary}

The intended soundness statement has a local layer and a physical layer.
The local layer says that if $\mathcal{K};\Gamma;\Delta\vdash t:\Gamma';\Delta'!E$ under a strict policy, then every acquisition atom in $t$ is justified by an allowed rule, every linear resource is used consistently with $\Delta$, and the reported effect is a conservative composition of rule effects.
This layer is represented today by Lean, Z3, Python checker obligations, and geometry checks.

The physical layer is conditional on the imported theorem dependencies attached to \Certified{} rules.
For example, the Steane/tetrahedral code-switch certificate locally checks the GF(2) algebraic gauge-fixing core and records the cited flag-circuit single-fault-tolerance theorem as an imported dependency.
Similarly, the Reed--Muller 15-to-1 certificate locally checks triorthogonality and low-weight $Z$-error detection, while the constant-time surface-code implementation profile is imported from an architecture-level result.
The paper's claim is therefore not that \system proves every physical gadget from first principles.
The claim is that, once local and imported rule contracts are made explicit, the calculus composes them without hiding illegal gates, duplicated resources, backend mismatches, or assumption boundaries.

\section{Implementation}

\lang is implemented as a small compiler/checker whose main job is to preserve the evidence that a fault-tolerant compilation step would otherwise erase.
The implementation is deliberately organized around three artifacts.
A source program is a list of \texttt{LogicalOp} nodes, which records what the algorithm asks for: preparation, logical gates, error correction, measurement, barriers, and classical branches.
A compiled target program is a tree of \texttt{FTStep} nodes, which records how each request was acquired: direct capability use, a switch path, a resource factory followed by consumption, or a derived region acquisition.
An \texttt{Effect} object travels with every step and gives the conservative account that the checker and the paper tables report.
This separation is the implementation analogue of the calculus judgment: source syntax keeps the user's logical story compact, while target syntax makes capability, resource, certificate, and backend evidence explicit.

The intermediate representation is small enough that the interesting invariants are visible in the data structures.
Each live logical qubit has a \texttt{QType(code, distance, mode)}, so a gate step does not merely name its operands; it also records the code in which those operands are currently protected.
Each target step stores \texttt{input\_codes}, \texttt{output\_codes}, an effect, a certificate level, and a free-form \texttt{details} map that carries rule-specific evidence such as source citations, layout events, produced resource identifiers, or consumed resource identifiers.
For example, a non-native $\gate{T}$ on \code{SurfaceD5} is not compiled into an opaque macro called ``apply T''.
It becomes a parent \texttt{acquire\_gate} step with two auditable children:
\[
  \mathsf{prepare\_resource}\ r:\res{TState};
  \quad
  \mathsf{consume\_resource\_gate}\ r\ \mathsf{for}\ \gate{T}.
\]
The identifier $r$ is part of the sidecar, not an explanatory comment, which is what allows the checker to reject a later attempt to consume the same resource again.

The planner treats acquisition as a local search over typed candidates rather than as a fixed gate-expansion pass.
When it encounters a gate $g$ on qubits currently in code $C$, it first checks whether the rule library contains a direct capability for $(C,g)$.
If not, it constructs candidates of two forms.
A switch candidate changes the operands to a code $C'$ that supports $g$, applies the direct rule there, and switches the operands back to $C$.
A resource candidate produces an appropriate resource state, such as a $\res{TState}$, $\res{BellPair}$, or $\res{CCZState}$, and consumes it for the requested operation.
The default score used by \texttt{best} mode is intentionally simple:
\[
  \mathsf{score}(p) =
  \mathsf{cycles}(p)
  + 0.01\,\mathsf{qubitRounds}(p)
  + 25\,\mathsf{switches}(p)
  + 35\,\mathsf{factories}(p)
  + 10^6\,\epsilon(p)
  + 20\,\mathsf{fail}(p).
\]
The point of this formula is not to claim an optimal cost model.
It gives the prototype a deterministic way to choose among qualitatively different acquisitions while leaving the chosen rule sequence, effect fields, and certificate assumptions available for review.

The hybrid strategy adds one compiler idea that is more useful than a gate-by-gate expansion in the case studies.
When the planner sees a contiguous block of gates and at least two of them need acquisition, it asks whether a single target code supports the whole block.
For the current rule library, \code{Tetra15} supports both $\gate{T}$ and $\gate{CNOT}$, while \code{SurfaceD5} supports neither directly.
The planner can therefore compile the block
\[
  \gate{CNOT};\ \gate{T};\ \gate{CNOT};\ \gate{T}
\]
as one \texttt{acquire\_region}: switch the involved qubits from \code{SurfaceD5} to \code{Tetra15}, run certified transversal gates, and switch them back.
The resulting target tree is useful precisely because it is not flattened into an anonymous cost.
The parent region records the source gate sequence and total effect; the children record each switch, each transversal gate, each layout event, and each certificate level.
Thus a reviewer can see that the transversal $\gate{T}$ and $\gate{CNOT}$ steps are certified through the \code{Tetra15} rules, while the compact \code{SurfaceD5}$\leftrightarrow$\code{Tetra15} bridge used in this prototype remains \Assumed{}.

Effect composition is implemented directly on the same object that is serialized into the sidecar.
Sequential composition adds logical-error upper bounds, cycles, qubit-rounds, decoder latency, gate counts, switch counts, and factory counts; it composes acceptance lower bounds multiplicatively; and it takes the maximum of peak-space fields.
Branch composition takes maxima for upper-bound quantities and minima for lower-bound quantities.
These operations are intentionally conservative, but their direct implementation has an important engineering advantage: the table entries in the paper, the SMT obligations emitted by the formal toolchain, and the checker's step summaries all refer to the same effect fields.
There is no separate presentation-only accounting layer.

The checker is a policy interpreter over the generated sidecar.
It flattens the target tree and checks five classes of obligations.
First, every code mentioned by a step must be supported by the selected backend topology.
Second, layout events must match the backend name, topology, capacity, free-qubit accounting, and switch workspace reservation reported by the layout component.
Third, direct \texttt{apply} steps must have a capability witness in the rule library.
Fourth, switch steps must name a known switch rule and that rule must be supported by the backend topology.
Fifth, resource identifiers must be produced before use and consumed at most once.
The certificate policy is checked at the same time: exploratory mode accepts \Assumed{} steps with warnings, whereas strict mode turns the same warnings into errors.
This policy is the implementation mechanism behind the paper's proof boundary.

The full-stack runner connects the compiler/checker to independent evidence generators.
It emits JSON sidecars for every case-study plan, SMT-LIB files for arithmetic bound checks, protocol certificates for selected GF(2) stabilizer and factory facts, decoder reports, geometry reports, artifact hashes, and a Lean build for the small calculus kernel.
The generated report currently records 53 obligations across code specifications, rule specifications, protocols, decoders, checker results, geometry, and Lean, together with 26 generated artifacts.
The report also carries an explicit certificate-boundary statement: GF(2), Z3, geometry, decoder-table checks, Python validation, and Lean discharge local obligations, while imported physical theorems and prototype assumptions remain visible dependencies.
This is the implementation invariant we rely on in the empirical sections: an accepted exploratory plan is auditable, but it is not silently promoted to a strict physics theorem.

\section{Case Study: Acquiring a Surface-Code Kernel}
\label{sec:case-study}

This case study follows a single \code{SurfaceD5} program through the planner and checker.
The program is small, but it is chosen to force the compiler to exercise the three acquisition mechanisms that matter for the language design: a $\gate{T}$ must be obtained from a magic-state resource or a code switch, a $\gate{CNOT}$ must be obtained from a Bell-style resource or a switch, and a $\gate{CCZ}$ must be obtained from a three-qubit resource state.
The source program in \texttt{experiments/run\_mvp.py} starts four logical qubits in \code{SurfaceD5}, runs a surface-code error-correction round, applies a native $\gate{H}$, enters a dense region containing two $\gate{CNOT}$ gates and two $\gate{T}$ gates, runs another correction round, applies an isolated $\gate{CCZ}$ and a final $\gate{T}$, and then uses a measured bit to guard a Clifford correction before the final observable measurement.
At the source level, this is just a logical kernel with gates, error correction, and classical control.
At the target level, it is a test of whether the compiler can explain every operation that \code{SurfaceD5} does not provide natively.

\begin{figure}[t]
\centering
\begin{tabular}{l}
\texttt{prepare q0:+; prepare q1:0; prepare q2:0; prepare r:0;}\\
\texttt{ec(q0, q1, q2, r, decoder=surface\_mwpm);}\\
\texttt{H(q0); barrier(dense\_acquisition\_region);}\\
\texttt{CNOT(q0,q1); T(q1); CNOT(q1,q2); T(q2);}\\
\texttt{ec(q0, q1, q2, decoder=surface\_mwpm);}\\
\texttt{barrier(isolated\_resource\_gate); CCZ(q0,q1,q2); T(r);}\\
\texttt{measure m0 = X(q0); if m0 then S(q2); measure out = Z(q2).}
\end{tabular}
\caption{The compact \code{SurfaceD5} kernel used in the standalone case study. The interesting part is not the logical algorithm itself, but the fact that \code{SurfaceD5} lacks direct $\gate{T}$, $\gate{CNOT}$, and $\gate{CCZ}$ capabilities in the rule library.}
\label{fig:surface-kernel}
\end{figure}

The first planner decision occurs at the dense acquisition barrier in Figure~\ref{fig:surface-kernel}.
If the planner works gate by gate, it can acquire each $\gate{CNOT}$ through the \texttt{BellTeleportCNOT\_SurfaceD5} resource rule and each $\gate{T}$ through the certified \texttt{Surface15to1TThenInject\_SurfaceD5} factory-and-injection rule.
This magic-oriented plan has a pleasing local structure: every unsupported operation leaves behind a fresh resource identifier, a factory effect, and a consumption step.
For the two $\gate{T}$ gates, the sidecar records the instantiated 15-to-1 profile with error $3.5\cdot 10^{-8}$, acceptance $0.985$, 6 cycles, 2775 qubit-rounds, and one factory each.
For the two $\gate{CNOT}$ gates, the current Bell-mediated profile is marked \Assumed{}, because the exact surface-code resource profile has not yet been imported from a specific physical protocol.
Exploratory mode keeps this plan available with warnings; strict mode rejects it.

The hybrid planner reads the same dense block differently.
Because \code{Tetra15} supports both transversal $\gate{CNOT}$ and transversal $\gate{T}$, the block
\[
  \gate{CNOT}(q_0,q_1);\quad
  \gate{T}(q_1);\quad
  \gate{CNOT}(q_1,q_2);\quad
  \gate{T}(q_2)
\]
can be acquired by changing the code of the participating data once, using the target code's native capabilities, and changing the data back.
The generated parent step is \texttt{HybridRegionAcquire\_Tetra15}.
Its children switch $q_0,q_1,q_2$ from \code{SurfaceD5} to \code{Tetra15}, apply certified transversal $\gate{CNOT}$, $\gate{T}$, $\gate{CNOT}$, and $\gate{T}$ rules, and switch the same qubits back to \code{SurfaceD5}.
The sidecar reports this region as 16 cycles, two switches, zero factories, error $6.5\cdot 10^{-5}$, and acceptance $0.941480149401$.
Those numbers make the compiler choice legible: the region avoids four separate factories or Bell acquisitions, but it pays for two prototype bridge switches whose assumptions remain explicit.

The layout sidecar explains a second reason why the region is a useful test.
Before the switch, the three surface-code patches occupy 75 physical qubits on the \texttt{Grid2D\_SurfaceLike} backend.
Switching them to \code{Tetra15} changes the live footprint to 45 qubits and records a workspace reservation of 93 qubits.
The event also records that the backend is \texttt{Grid2D\_SurfaceLike} and the topology is \texttt{grid2d}.
The checker does not prove a detailed lattice-surgery route from this event, but it does check the conservative accounting: the event names the selected backend, uses a supported topology, does not exceed capacity, and has consistent free-qubit arithmetic.
This is the level of layout evidence the current prototype claims.

After the dense block, the program returns to \code{SurfaceD5} and runs another error-correction round.
That transition is important for the story because it demonstrates that code-indexed typing is not merely a way to choose a gate implementation.
The checker sees the output codes of the hybrid region as \code{SurfaceD5}, so the subsequent \texttt{EC\_SurfaceD5} rule and final \texttt{MeasureL\_SurfaceD5} rules are applicable.
Had the region failed to switch a qubit back, the later surface-code rules would no longer match the type environment.
The type index therefore records a real protocol state that later steps depend on.

The isolated $\gate{CCZ}$ gate exercises a different path.
There is no direct \code{SurfaceD5} capability for $\gate{CCZ}$ and the hybrid region has already ended, so the planner emits a resource acquisition using \texttt{CCZFactoryThenInject\_SurfaceD5}.
The sidecar first produces \texttt{res\_ccz\_1} as a $\res{CCZState}$ and then consumes exactly \texttt{res\_ccz\_1} for $\gate{CCZ}(q_0,q_1,q_2)$.
The current profile contributes 34 cycles, one factory, 5200 qubit-rounds, acceptance $0.88$, and an \Assumed{} certificate level.
This is intentionally written as a placeholder in the rule library: it keeps the case study multi-qubit and tests the resource-linearity path, but it is not presented as a paper-backed CCZ factory theorem.
Strict mode therefore rejects the plan until that placeholder is replaced by a certified rule.

The final $\gate{T}$ on the qubit $r$ shows how a certified resource rule appears beside assumed prototype rules in the same plan.
The planner chooses the high-fidelity surface-code 15-to-1 factory rule \texttt{Surface15to1TThenInject\_SurfaceD5}.
The generated children produce \texttt{res\_t\_4} and then consume it for $\gate{T}(r)$.
The rule is marked \Certified{} and records the instantiated profile at distance 5: output error $35p^3=3.5\cdot 10^{-8}$ for $p=10^{-3}$, 6 cycles, $111d^2=2775$ qubit-rounds, acceptance approximated as $1-15p=0.985$, and peak workspace 775.
The implementation does not locally reprove the physical surface-code schedule behind these constants; it records that schedule as an imported theorem dependency and checks the local resource and accounting structure around it.

The checker result is best understood as a certificate audit rather than a yes/no benchmark.
For the hybrid plan, exploratory checking succeeds and reports warnings for the \code{SurfaceD5}$\leftrightarrow$\code{Tetra15} bridge and the placeholder $\res{CCZState}$ factory.
Strict checking fails because those warnings correspond to \Assumed{} rules.
For the best-scoring plan under the current cost formula, the planner selects resource acquisitions for the two $\gate{CNOT}$ gates, the two dense-block $\gate{T}$ gates, the $\gate{CCZ}$, and the final $\gate{T}$; it reports 90 cycles, error bound $2.56125\cdot 10^{-4}$, acceptance lower bound $0.83763042984412$, six factories, and no switches.
The hybrid plan instead reports 76 cycles, error bound $2.71055\cdot 10^{-4}$, acceptance lower bound $0.8160749935007868$, two factories, and two switches.
Both plans are useful compiler artifacts in exploratory mode.
Neither is counted as a strict physical certification because both still contain assumed prototype components.

The case study therefore illustrates the central claim of \system.
A logical program can remain simple, but the generated artifact must not be simple in the same way.
It must remember why a non-native gate is legal, where a resource came from, where that resource was consumed, how code indices changed across a region, how conservative costs were composed, and which physical facts were locally checked, imported, or assumed.
The \code{SurfaceD5} kernel is not meant to be a performance benchmark for fault-tolerant architectures.
It is a compact end-to-end demonstration that the language can turn fault-tolerant acquisition choices into checkable program structure.

\section{Evaluation}

The evaluation asks whether the prototype produces auditable acquisition evidence for the kinds of plans described in Section~\ref{sec:case-study}.
We evaluate the planner/checker on two logical workloads over the same \texttt{Grid2D\_SurfaceLike} backend.
The first is the compact \code{SurfaceD5} kernel from the case study.
The second is a denser phase-oracle workload with seven initial surface-code qubits, a longer region of $\gate{CNOT}$ and $\gate{T}$ operations, two measured branches, a $\gate{CCZ}$ acquisition, and final logical measurements.
The goal is not to beat a production FTQC compiler.
It is to check that different acquisition strategies produce different sidecars, that the checker accepts those sidecars only under the right certificate policy, and that negative examples fail for the intended static reasons.

\begin{table}[t]
\centering
\caption{Compact \code{SurfaceD5} kernel on \texttt{Grid2D\_SurfaceLike}. All rows are accepted in exploratory mode and rejected in strict mode because each contains at least one \Assumed{} rule. Lower is better for error, cycles, qubit-rounds, switches, and factories; higher is better for acceptance.}
\label{tab:surface-kernel-results}
\begin{tabular}{lrrrrrr}
\toprule
Strategy & Err. & Accept & Cycles & Q-rounds & Switch & Factory \\
\midrule
magic-only & $2.56125{\cdot}10^{-4}$ & 0.8376 & 90 & 23565 & 0 & 6 \\
switch-only & $3.3302{\cdot}10^{-4}$ & 0.7959 & 111 & 27313 & 10 & 1 \\
hybrid & $2.71055{\cdot}10^{-4}$ & 0.8161 & 76 & 23705 & 2 & 2 \\
best & $2.56125{\cdot}10^{-4}$ & 0.8376 & 90 & 23565 & 0 & 6 \\
random seed 7 & $2.98055{\cdot}10^{-4}$ & 0.8293 & 104 & 24045 & 4 & 4 \\
random seed 17 & $4.56055{\cdot}10^{-4}$ & 0.4230 & 128 & 20015 & 0 & 6 \\
\bottomrule
\end{tabular}
\end{table}

Table~\ref{tab:surface-kernel-results} summarizes the compact kernel without repeating the full case-study trace.
The main observation is that the sidecar exposes a real strategy tradeoff.
The magic-only and best rows coincide under the current score because the resource path keeps the error and acceptance accounting favorable despite using six factories.
The hybrid row is faster, at 76 cycles, because it acquires the two $\gate{CNOT}$ gates and two $\gate{T}$ gates in the dense block through a single \code{Tetra15} region, but it pays two switch events and records bridge assumptions.
The switch-only row has fewer factories but ten switches and a worse conservative error/acceptance summary.
The randomized rows are useful sanity checks: they show that the planner is not hard-wired to one expansion and that mixed acquisition choices produce distinct cost and certificate profiles.

\begin{table}[t]
\centering
\caption{Denser phase-oracle workload on the same backend. The hybrid strategy reduces cycle count by acquiring the first dense non-native region once, while the best score still prefers the magic-only row under the current weighted objective.}
\label{tab:complex-results-sections}
\begin{tabular}{lrrrrrr}
\toprule
Strategy & Err. & Accept & Cycles & Q-rounds & Switch & Factory \\
\midrule
magic-only & $5.12195{\cdot}10^{-4}$ & 0.7957 & 147 & 42630 & 0 & 12 \\
switch-only & $6.6802{\cdot}10^{-4}$ & 0.7056 & 189 & 51640 & 22 & 1 \\
hybrid & $5.42055{\cdot}10^{-4}$ & 0.7553 & 107 & 42910 & 2 & 4 \\
best & $5.12195{\cdot}10^{-4}$ & 0.7957 & 147 & 42630 & 0 & 12 \\
random seed 23 & $8.5159{\cdot}10^{-4}$ & 0.2686 & 215 & 43389 & 6 & 9 \\
\bottomrule
\end{tabular}
\end{table}

The denser workload in Table~\ref{tab:complex-results-sections} stresses the same mechanism with more opportunities to make a poor acquisition choice.
The magic-only plan uses twelve factories and no switches, yielding 147 cycles and acceptance 0.7957.
The switch-only plan uses one factory but twenty-two switches, yielding 189 cycles and acceptance 0.7056.
The hybrid plan uses two switches and four factories, yielding 107 cycles and acceptance 0.7553.
This is the empirical pattern the case study was designed to make understandable: a region switch can reduce time when several non-native gates share a target code, but the checker still records the price paid in switches, assumptions, and conservative acceptance.
Because the objective function also weights error, qubit-rounds, switches, and factories, \texttt{best} currently selects the magic-only row rather than the fastest row.
That behavior is acceptable for this prototype because the objective is transparent and the generated sidecars allow the paper to report the tradeoff instead of hiding it behind a single winner.

The negative tests in Table~\ref{tab:negative-tests-sections} validate the failure modes that motivate the calculus.
The first test constructs a target step that directly applies $\gate{CNOT}$ to two \code{SurfaceD5} qubits.
The checker rejects it twice, once for each operand, because the rule library contains no direct $\Cap(\code{SurfaceD5},\gate{CNOT})$ witness.
The second test produces one fake $\res{TState}$ identifier and consumes it for two different target gates.
The checker rejects the plan with \texttt{resource res\_t\_shared consumed twice}, which is the implementation analogue of removing a resource from the linear context after first use.
The third test attempts to switch from \code{SurfaceD5} to the negative-control \code{QLDPC12} profile on the grid backend.
The checker rejects both the target code and the switch rule because the profile requires long-range connectivity; exploratory mode also reports that the switch rule is \Assumed{}.
These tests are small by design, but each corresponds to a concrete static obligation in the formal calculus.

\begin{table}[t]
\centering
\caption{Negative tests. Each rejected plan targets a different checker invariant.}
\label{tab:negative-tests-sections}
\begin{tabular}{lp{0.52\linewidth}p{0.22\linewidth}}
\toprule
Test & Rejected behavior & Checker reason \\
\midrule
direct CNOT & Apply $\gate{CNOT}$ directly to \code{SurfaceD5} data. & Missing direct capability witness. \\
duplicate resource & Consume the same \texttt{res\_t\_shared} identifier twice. & Linear resource already consumed. \\
qLDPC switch & Switch \code{SurfaceD5} to \code{QLDPC12} on a grid backend. & Code and rule require long-range topology. \\
\bottomrule
\end{tabular}
\end{table}

The integrated evidence report checks that these examples are not isolated scripts.
The current full-stack run records 11 case-study compile results, 11 geometry reports, 3 code-specification groups, 2 rule-specification groups, 2 protocol certificates, 2 decoder reports, 53 total obligations, and 26 generated artifacts with hashes.
The protocol layer checks local GF(2) facts for selected switch and factory certificates; the decoder layer distinguishes a checked small-code syndrome table from an imported surface-code decoder contract; the geometry layer checks conservative rectangular embeddings and workspace accounting; the SMT files encode arithmetic bound obligations; and the Lean build checks the small resource/certificate core.
The report's boundary statement is part of the result: locally discharged obligations and imported physical theorem dependencies are reported separately, and \Assumed{} prototype rules are never counted as local proofs.

The evaluation has three limitations.
First, the workloads are case studies rather than a benchmark suite, so the numbers should be read as evidence about the compiler/checker interface, not as architecture-level performance claims.
Second, several useful bridge and resource rules are still \Assumed{}, which means the current case-study plans are exploratory artifacts and strict mode correctly rejects them.
Third, the layout checker uses conservative capacity and rectangular-embedding checks rather than a complete scheduler for lattice surgery, routing, or factory placement.
These limitations are not peripheral: they define the current proof boundary.
The empirical claim supported by the present evaluation is that \system can generate, check, and report auditable acquisition plans with explicit costs and certificate levels, including clear rejection of illegal capability use, duplicate resources, and backend-incompatible switches.

\section{Related Work}

\paragraph{Quantum programming languages and substructural types.}
Quantum programming languages use linear and affine type disciplines to reflect no-cloning and controlled disposal of quantum data.
Recent POPL work on Qurts uses Rust-like lifetimes to control automatic uncomputation through a substructural discipline~\cite{hirata2025qurts}.
\system shares the view that quantum programming needs static ownership information, but its ownership discipline applies to encoded code states and fault-tolerant resources rather than to automatic uncomputation.
In particular, a $\res{TState}$ in \system is treated as a linear physical artifact whose consumption is part of the target acquisition plan.

\paragraph{Resource estimation for quantum circuit descriptions.}
Type-based resource estimation for quantum circuit description languages uses effects, refinements, and index terms to derive circuit-size bounds~\cite{colledan2025resource}.
\system also uses effects, but the quantities are specialized to fault-tolerant acquisition: logical error, failure, acceptance, code switches, factories, decoder latency, physical qubits, cycles, and qubit-rounds.
The effect system is therefore tied to rule certificates and backend side conditions, not only to circuit size.
This distinction is visible in the case study, where two plans can implement the same logical gate sequence but expose different switch/factory/certificate profiles.

\paragraph{Quantum compiler generation and mapping.}
Compiler frameworks for qubit mapping and routing identify common state-machine structure across diverse architectures and synthesize compilers from architecture specifications~\cite{molavi2026qmr}.
\system is complementary.
It does not attempt to solve low-level mapping and routing optimally; instead, it exposes code capabilities and fault-tolerant acquisition decisions as typed program structure.
The current geometry checker is a conservative rectangular embedding and workspace-capacity check, not a full nearest-neighbor routing or lattice-surgery scheduler.

\paragraph{Parameterized quantum verification.}
Automata-based verification can prove relational properties of parameterized quantum circuits across all input sizes~\cite{abdulla2026parameterized}.
\system targets a different verification layer.
Rather than verifying arbitrary circuit families against functional specifications, it checks that a compiler-generated fault-tolerant acquisition plan respects code capabilities, linear resources, certificate levels, backend constraints, and conservative effect bounds.
The two approaches are compatible: parameterized circuit verification can supply facts about circuit families, while \system can record how such facts are used in an acquisition plan.

\paragraph{Fault-tolerant protocols.}
\system relies on the standard magic-state view of universal fault-tolerant computation~\cite{bravyi2005magic} and recent protocol results for code switching, magic-state distillation, and surface-code implementation profiles~\cite{butt2024switching,heussen2025oneway,itogawa2025zero,wan2024constant}.
The calculus does not replace those protocol proofs.
It provides a language-level composition layer that records which protocol facts are locally checked, which are imported, and which are still assumptions in the prototype rule library.

\section{Limitations and Future Work}

The current draft has a deliberately narrow soundness boundary.
The Lean development mechanizes a small core of capability, certificate, resource, and effect properties, but not yet a full source-to-target elaboration theorem.
The Python checker validates generated plans against the current sidecar format, but the implementation does not yet extract a complete machine-checkable typing derivation for every accepted plan.
Closing this gap is the most important next step for a POPL submission.

Several rules in the current case studies are marked \Assumed{}.
These rules are useful because they let the compiler exercise end-to-end acquisition strategies that connect otherwise paper-backed pieces, but they are not physics claims.
The paper therefore reports exploratory results with warnings and treats strict mode as the proof-facing policy.
Before submission, prototype bridge rules should either be replaced by literature-backed certificates or moved into a clearly separated exploratory appendix.

The layout model is also conservative.
The geometry checker verifies final rectangular patch embeddings and workspace reservations on a fixed backend grid, but it does not model complete patch routing, lattice-surgery boundary compatibility, or factory scheduling optimality.
This limitation is acceptable for the current language claim because topology and workspace violations are already visible to the checker.
It is not sufficient for claiming a production-quality FTQC layout compiler.

The benchmark set is small.
The case studies validate the calculus and checker interface, but they do not compare against production FTQC compilers or a broad suite of fault-tolerant workloads.
Future work should add larger logical programs, more backend profiles, stronger acquisition baselines, and ablations that isolate the effect of hybrid-region planning, certificate policies, and resource-factory choices.

\section{Conclusion}

\system turns fault-tolerant acquisition decisions into typed, auditable program structure.
The key idea is to keep the code index, linear resource context, conservative effect summary, backend side conditions, and certificate level visible while a logical program is elaborated into an executable acquisition plan.
In the running surface-code example, this visibility is what distinguishes a legal direct $\gate{H}$, an acquired $\gate{T}$, a hybrid switched region, and an illegal direct $\gate{CNOT}$.

The current implementation demonstrates the interface end to end across planning, checking, SMT obligations, protocol certificates, decoder contracts, geometry checks, and a Lean proof kernel.
The evidence is intentionally conditional: local algebraic and software-level obligations are checked locally, while complete physical protocol claims remain imported theorem dependencies.
This boundary is the paper's main programming-language message.
Certified local FTQC rules can be composed through a small, code-indexed calculus that makes legality, resources, effects, layout constraints, and proof assumptions explicit.


\begin{acks}
Omitted for anonymous review.
\end{acks}

\bibliographystyle{ACM-Reference-Format}
\bibliography{references}

\appendix
\section{Draft Claim--Evidence Map}

\begin{table}[h]
\centering
\caption{Major claims in the current draft and their supporting artifacts.}
\begin{tabular}{p{0.32\linewidth}p{0.45\linewidth}p{0.15\linewidth}}
\toprule
Claim & Evidence & Status \\
\midrule
Code-indexed acquisition plans expose legality of native gates, switches, resources, and backends.
& \texttt{cipr/ir.py}, \texttt{cipr/planner.py}, \texttt{cipr/checker.py}, \texttt{cipr/rules.py}. & Supported \\
Linear resources prevent duplicate consumption.
& Python negative test; Lean theorems \texttt{no\_double\_consume\_same\_resource} and \texttt{badDoubleConsumePlan\_not\_linear\_after\_single\_resource}. & Supported \\
Strict policy rejects assumed prototype rules.
& Python checker \texttt{allow\_assumed=False}; Lean theorem \texttt{strict\_rejects\_assumed\_surface\_to\_tetra}. & Supported \\
Protocol-level physical proofs are explicit imported dependencies.
& \texttt{experiments/outputs/full\_stack\_report.json}; protocol certificate records for Butt 2024 and Wan 2024. & Supported \\
End-to-end compiler stack is reproducible.
& Repository commands in \texttt{README.md}; full-stack report; current local evidence from prior run. & Supported, requires rerun before submission \\
\bottomrule
\end{tabular}
\end{table}


\end{document}
